\par\medskip{\centering\scriptsize\begin{tabular}{lrrrrrrr}\toprule
 & \multicolumn{4}{c}{Dette 2047 (points)} & \multicolumn{3}{c}{Intérêts 2047 (points de PIB)} \\\cmidrule(lr){2-5}\cmidrule(lr){6-8}
 & sans M11 & prime 0 & 0,03 & 0,05 & prime 0 & 0,03 & 0,05 \\ \midrule
S1 Protection & $+4$ & $+8$ & $+10$ & $+11$ & $+0{,}6$ & $+1{,}0$ & $+1{,}3$ \\
S2 Emploi & $+15$ & $+18$ & $+26$ & $+33$ & $+1{,}2$ & $+2{,}7$ & $+4{,}3$ \\
S3 Productivité & $+34$ & $+37$ & $+53$ & $+69$ & $+2{,}0$ & $+5{,}4$ & $+9{,}0$ \\
S4 Redressement & $-43$ & $-44$ & $-59$ & $-72$ & $-2{,}0$ & $-4{,}9$ & $-7{,}4$ \\
S5 Global & $+2$ & $+3$ & $+5$ & $+7$ & $+0{,}3$ & $+0{,}7$ & $+1{,}1$ \\
\bottomrule\end{tabular}\par}\medskip
 La réaction de la BCE et le commerce avec la zone euro comptent peu (colonne « prime 0 ») : 1 à 4 points de dette. L'essentiel vient de la \emph{prime de risque} sur les taux français, de 0,03 point de taux long par point de dette (hypothèse de niveau C, au milieu des estimations de la littérature, 0,02 à 0,05) : un écart de dette se renchérit ou s'allège par les intérêts, et l'effet se cumule sur vingt ans. Pour S3, la dette de 2047 passe de $+$37 points sans prime à $+$53 avec la prime centrale et $+$69 avec une prime de 0,05. Cette hypothèse n'était pas variée dans les configurations faible et forte ; elle est désormais traitée comme une incertitude structurelle à part entière.
